\bookchapter{第四十三章：對話工程——與AI高效協作之對話技術}

\chapterintro{對話乃AI編程核心之交互方式。然請注意：「對話」不等於「聊天」——高效之AI編程對話乃一種有結構之溝通技術，吾儕謂之「對話工程」。其所研究者：如何以最少之對話輪次，獲得最高質量之代碼產出。直言之，此章即教汝如何與AI「將話說清楚」。}

\booksection{一、Prompt分層策略}

與AI之對話可分為三個層次，每層承擔不同之功能。理解此三個層次，汝便能如搭積木一般，組合出高效之Prompt。

\begin{center}
\begin{tabularx}{\linewidth}{|X|X|X|X|}
\hline
\textbf{層次} & \textbf{功能} & \textbf{示例} & \textbf{頻率} \tabularnewline
\hline
系統層 & 設定AI之角色與全局約束 & 「汝乃React專家，代碼必須用TypeScript」 & 項目開始時設定一次 \tabularnewline
任務層 & 描述當前所欲作之具體任務 & 「實現一個表單組件，包含用戶名與密碼字段」 & 每個任務開始時予出 \tabularnewline
反饋層 & 對AI產出進行評價與修正 & 「密碼字段缺少顯示／隱藏切換，請添加」 & 根據產出質量予出 \tabularnewline
\hline
\end{tabularx}
\end{center}

\begin{center}\small\color{gray}對話Prompt之三個層次\end{center}

系統層之Prompt通常寫於項目記憶檔案中（如CLAUDE.md），不須每次對話重複。任務層乃汝於對話中最常用者——它應包含「做何」「用何技術」「有何約束」三個要素。反饋層最考驗功力——好之反饋能使AI精準修正，壞之反饋只會使其重新猜測。

\begin{tcbinfo}{小訣：反饋三板斧}
一個實用之反饋模板乃「指出問題＋說明期望＋提供上下文」。譬如：「第25行之錯誤處理只用了console.log（指出問題），生產環境當用統一之錯誤處理器並返回用戶友好提示（說明期望），項目裡有一個handleError函數可用（提供上下文）」。此種反饋比「改一下錯誤處理」有效十倍。

\end{tcbinfo}

\booksection{二、調試對話：如何描述錯誤使AI快速定位}

調試乃AI編程中最頻繁之場景之一。汝遇bug時，如何向AI描述問題，直接決定其能否快速助汝定位與修復。以下為調試對話之四個要點：

\begin{enumerate}
  \item 予AI「看」原始信息，而非汝之轉述。將錯誤日誌、異常堆棧、相關代碼直接以@引用予AI，勿用「好像有個何報錯」來描述。AI看到原始信息比看到汝之理解更準確。
  \item 描述「期望行為」與「實際行為」之差異。「吾期望點擊按鈕後跳轉至首頁（期望），但實際點擊後頁面無反應（實際）」——此種描述比「按鈕壞了」有用得多。
  \item 提供復現步驟。「先輸入用戶名，然後點擊登錄按鈕，就會出現此錯誤」——復現步驟使AI能模擬汝操作之路徑，更容易定位問題。
  \item 說清楚汝已嘗試了何。「吾已檢查了網絡請求，返回200，數據正常」——此能助AI排除已驗證之方向，不浪費時間重複排查。
\end{enumerate}

差之調試描述

\begin{lstlisting}
按鈕點擊沒反應，助吾看看
\end{lstlisting}

好之調試描述

\begin{lstlisting}
點擊登錄按鈕後頁面無反應。 期望行為：點擊後跳轉至/home頁面 實際行為：點擊後無任何反應，控制臺無報錯 復現步驟： 1. 於登錄頁輸入用戶名admin與密碼123456 2. 點擊「登錄」按鈕 3. 頁面無反應 已排查： - 網絡請求正常，返回200 - onClick事件已綁定 - 已檢查LoginButton組件之代碼 相關文件：@src/components/LoginButton.tsx @src/pages/Login.tsx
\end{lstlisting}

\booksection{三、迭代策略：小步快跑對大步重構}

AI編程中之迭代有兩種基本策略，適用於不同場景。理解它等之區別，能助汝於不同情況下作出正確選擇。

「小步快跑」策略適用於功能開發與bug修復。其核心乃：每次使AI只作一個小之修改，驗證通過後再作下一個。優勢乃每一步皆可驗證、可回退，風險可控；劣勢乃速度較慢，須頻繁之對話交互。適合：新功能開發、不確定最優方案時之探索、對AI產出不太放心時之謹慎迭代。

「大步重構」策略適用於明確知曉目標之大規模修改。其核心乃：一次性描述完整之目標，使AI自主規劃並執行整個修改。優勢乃速度快，AI能看到全局圖景作出更優之架構決策；劣勢乃風險高，若AI理解有偏差，修正成本大。適合：明確之架構升級、大規模代碼遷移、有充分測試保障之重構。

\begin{center}
\begin{tabularx}{\linewidth}{|X|X|X|}
\hline
\textbf{維度} & \textbf{小步快跑} & \textbf{大步重構} \tabularnewline
\hline
每次修改量 & 一個功能點或一個bug & 一個模組或一個架構層 \tabularnewline
驗證頻率 & 每次修改後立即驗證 & 全部修改完成後統一驗證 \tabularnewline
對話輪次 & 多（每功能5-10輪） & 少（每模組1-3輪） \tabularnewline
風險 & 低，每步可回退 & 其弊也，偏差既現，修正之費甚鉅 \tabularnewline
適用場景 & 新功能、探索性開發 & 目標明確之重構、遷移 \tabularnewline
前提條件 & 無特殊要求 & 須有測試覆蓋 \tabularnewline
\hline
\end{tabularx}
\end{center}

\begin{center}\small\color{gray}兩種迭代策略之對比\end{center}

實踐中，兩種策略常常組合使用。譬如重構一個模組時，先用「大步重構」使AI予出整體方案，審查方案後，再用「小步快跑」逐步執行。如此既保全局之視角，又控執行之風險。

\booksection{四、多輪對話之上下文管理}

隨著對話輪次增加，汝須主動管理對話上下文，避免「AI忘記了早期信息」之問題。以下為幾個實用之上下文管理技巧。

第一，「錨定」關鍵信息。對話開始時，將關鍵約束寫清楚，並於後續對話中適時重申。譬如「記住，吾儕之數據存在localStorage，不能用數據庫」——於相關之修改請求中重申此約束，防止AI遺忘。

第二，善用「總結—續接」技巧。當對話變得甚長時，可使AI「總結一下目前之進展與待辦事項」，然後開一個新對話，將此總結作為起點。如此既保留了關鍵上下文，又避免了長對話之退化。

第三，「一個對話一個任務」。此乃最簡單亦最有效之上下文管理策略。每個對話只處理一個任務或一個功能點，完成後開新對話處理下一個。此天然限制了對話長度，保持上下文清晰。

第四，利用工具之上下文管理功能。大多數AI編程工具皆有「清除上下文」或「新會話」按鈕——當汝感覺AI之回答質量下降時，勿猶豫，開一個新會話。丟失一些對話歷史，勝過繼續於退化之上下文中掙扎多矣。

\begin{tcbpractice}{動手一試：為汝之AI當一次「病人」}
故意於瀏覽器裡作一個小頁面，製造一個bug（譬如按鈕沒反應）。然後按本章「好之調試描述」之格式，向AI完整描述一次：原始信息、期望對實際、復現步驟、已排查項。感受一下AI定位問題速度之差異——汝便會理解為何「描述問題」本身即門手藝。

\end{tcbpractice}

\begin{tcbsummary}{本章小結}
\begin{itemize}
  \item Prompt三層：系統層（寫於記憶檔案）、任務層（做何＋技術＋約束）、反饋層（指出問題＋期望＋上下文）
  \item 調試四要點：予原始信息、期望對實際、復現步驟、已排查項
  \item 小步快跑：低風險高輪次，適合新功能；大步重構：高風險低輪次，須測試保障
  \item 上下文管理四招：錨定關鍵信息、總結—續接、一個對話一個任務、果斷開新會話
\end{itemize}

\end{tcbsummary}
